\documentclass[12pt,a4paper]{article}
\usepackage{amsmath, amssymb}
\usepackage[margin=2.5cm]{geometry}
\usepackage{enumitem}
\usepackage{xcolor}

\newcommand{\easy}{\textcolor{green!60!black}{$\bigstar$}}
\newcommand{\medium}{\textcolor{orange}{$\bigstar\bigstar$}}
\newcommand{\hard}{\textcolor{red}{$\bigstar\bigstar\bigstar$}}

\title{\textbf{Solutions: Additional Exercises 2}\\General Series}
\author{Series and Integral Transforms}
\date{}

\begin{document}
\maketitle

% -------------------------------------------------------
\section*{Exercise 1 — Telescoping Sums (Exact Values)}

\subsection*{(a) \easy\ $\displaystyle\sum_{n=1}^{\infty}\frac{1}{n(n+1)}$}

Partial fractions: $\dfrac{1}{n(n+1)}=\dfrac{1}{n}-\dfrac{1}{n+1}$.  Telescoping:
\[
  S_N = 1 - \frac{1}{N+1}\;\to\; \boxed{1}.
\]

\subsection*{(b) \easy\ $\displaystyle\sum_{n=1}^{\infty}\frac{1}{n(n+2)}$}

Partial fractions: $\dfrac{1}{n(n+2)}=\dfrac{1}{2}\!\left(\dfrac{1}{n}-\dfrac{1}{n+2}\right)$.

This is a \emph{step-2 telescope}. Grouping:
\[
  S_N = \frac{1}{2}\!\left(1+\frac{1}{2}-\frac{1}{N+1}-\frac{1}{N+2}\right)
  \;\to\; \frac{1}{2}\cdot\frac{3}{2} = \boxed{\frac{3}{4}}.
\]

\subsection*{(c) \medium\ $\displaystyle\sum_{n=2}^{\infty}\ln\!\left(1-\frac{1}{n^2}\right)$}

Factor: $1-\dfrac{1}{n^2}=\dfrac{n^2-1}{n^2}=\dfrac{(n-1)(n+1)}{n^2}$.  Using log rules:
\[
  \ln\!\left(\frac{(n-1)(n+1)}{n^2}\right)
  = \underbrace{\bigl(\ln(n-1)-\ln n\bigr)}_{u_n}
  + \underbrace{\bigl(\ln(n+1)-\ln n\bigr)}_{v_n}.
\]
For the partial sum $S_N$:
\begin{align*}
  \sum_{n=2}^{N} u_n &= (\ln 1-\ln 2)+(\ln 2-\ln 3)+\cdots+(\ln(N-1)-\ln N)
    = \ln 1-\ln N = -\ln N,\\
  \sum_{n=2}^{N} v_n &= (\ln 3-\ln 2)+(\ln 4-\ln 3)+\cdots+(\ln(N+1)-\ln N)
    = \ln(N+1)-\ln 2.
\end{align*}
Hence $S_N = -\ln N + \ln(N+1) - \ln 2 = \ln\!\left(\dfrac{N+1}{N}\right) - \ln 2 \to 0 - \ln 2 = \boxed{-\ln 2}$.

\subsection*{(d) \medium\ $\displaystyle\sum_{n=1}^{\infty}\frac{1}{n(n+1)(n+2)}$}

Use the identity $\dfrac{1}{n(n+1)(n+2)}=\dfrac{1}{2}\!\left(\dfrac{1}{n(n+1)}-\dfrac{1}{(n+1)(n+2)}\right)$.

Telescoping:
\[
  S_N = \frac{1}{2}\!\left(\frac{1}{1\cdot2}-\frac{1}{(N+1)(N+2)}\right)
  \;\to\; \frac{1}{2}\cdot\frac{1}{2} = \boxed{\frac{1}{4}}.
\]

% -------------------------------------------------------
\section*{Exercise 2 — Alternating Series: Classification}

For each series we determine whether it \textbf{diverges}, is \textbf{conditionally} convergent, or \textbf{absolutely} convergent.

\subsection*{(a) \easy\ $\displaystyle\sum\frac{(-1)^{n+1}}{n^2+1}$}

$|a_n|=\dfrac{1}{n^2+1}\leq\dfrac{1}{n^2}$ and $\sum 1/n^2$ converges ($p=2$). By Comparison, $\sum|a_n|$ converges.\\
\textbf{Absolutely convergent}.

\subsection*{(b) \easy\ $\displaystyle\sum\frac{(-1)^n}{\sqrt{n}}$}

$\sum 1/\sqrt{n}$ diverges ($p=1/2$), so not absolutely convergent.  However, $1/\sqrt{n}\searrow 0$, so \textbf{Leibniz's Test} applies.\\
\textbf{Conditionally convergent}.

\subsection*{(c) \easy\ $\displaystyle\sum(-1)^n\frac{n+1}{n}$}

$a_n = \dfrac{n+1}{n}\to 1\neq 0$. By the \textbf{Necessary Condition} the series \textbf{diverges}.

\subsection*{(d) \medium\ $\displaystyle\sum(-1)^n\frac{\ln n}{n}$, $n\geq 2$}

$\sum\dfrac{\ln n}{n}$: since $\ln n\geq 1$ for $n\geq 3$, we have $\dfrac{\ln n}{n}\geq\dfrac{1}{n}$, so by Comparison with the harmonic series, $\sum|\,a_n|$ \textbf{diverges}.

For conditional convergence: $f(x)=\ln x/x$ satisfies $f'(x)=(1-\ln x)/x^2 <0$ for $x>e$, so the sequence is eventually decreasing, and $(\ln n)/n\to 0$.  By \textbf{Leibniz's Test}, the series \textbf{converges conditionally}.

\subsection*{(e) \easy\ $\displaystyle\sum\frac{(-1)^n}{n(n+1)}$}

$\dfrac{1}{n(n+1)}\sim\dfrac{1}{n^2}$; \textbf{Limit Comparison} with $1/n^2$ gives $L=1$, so $\sum|a_n|$ converges.\\
\textbf{Absolutely convergent}.

\subsection*{(f) \medium\ $\displaystyle\sum(-1)^n\!\left(1-\cos\frac{1}{n}\right)$}

For small $u$, $1-\cos u\sim u^2/2$, so $1-\cos(1/n)\sim 1/(2n^2)$.  \textbf{Limit Comparison}: $\lim n^2(1-\cos(1/n))=1/2$.  Since $\sum 1/n^2$ converges, $\sum|a_n|$ converges.  \textbf{Absolutely convergent}.

\subsection*{(g) \medium\ $\displaystyle\sum(-1)^n\!\left(\sqrt{n+1}-\sqrt{n}\right)$}

Rationalise: $\sqrt{n+1}-\sqrt{n}=\dfrac{1}{\sqrt{n+1}+\sqrt{n}}\sim\dfrac{1}{2\sqrt{n}}$.

Absolute series: $\sum\dfrac{1}{\sqrt{n+1}+\sqrt{n}}\sim\sum\dfrac{1}{2\sqrt{n}}$ diverges.\\
The terms decrease to $0$, so \textbf{Leibniz's Test} applies. \textbf{Conditionally convergent}.

\subsection*{(h) \medium\ $\displaystyle\sum(-1)^n\!\left(1-\!\left(\frac{n}{n+1}\right)^n\right)^{\!1/2}$}

Now $\left(\dfrac{n}{n+1}\right)^n = \left(1-\dfrac{1}{n+1}\right)^n\to e^{-1}$, so:
\[
  1 - \left(\frac{n}{n+1}\right)^n \to 1 - \frac{1}{e} \neq 0
  \implies a_n \to \sqrt{1-1/e}\neq 0.
\]
\textbf{Diverges} by the Necessary Condition.

\subsection*{(i) \hard\ $\displaystyle\sum(-1)^n\frac{(2n)!}{4^n(n!)^2}$}

By Stirling: $\dfrac{(2n)!}{4^n(n!)^2}\sim\dfrac{1}{\sqrt{\pi n}}\to 0$, and this is decreasing (ratio $(2n+1)/(2n+2)<1$).
\textbf{Leibniz's Test} applies $\Rightarrow$ converges.  But $\sum\dfrac{1}{\sqrt{\pi n}}$ diverges ($p=1/2$).\\
\textbf{Conditionally convergent}.

% -------------------------------------------------------
\section*{Exercise 3 — General Convergence}

\subsection*{(a) \easy\ $\displaystyle\sum\frac{\cos(2\pi n)}{n^2}$}

Note $\cos(2\pi n)=1$ for all $n\in\mathbb{N}$.  So the series equals $\sum 1/n^2$, which converges ($p=2$).  \textbf{Absolutely convergent}.

\subsection*{(b) \easy\ $\displaystyle\sum\frac{n}{e^n}$}

\textbf{Ratio Test}: $\dfrac{(n+1)/e^{n+1}}{n/e^n}=\dfrac{n+1}{en}\to\dfrac{1}{e}<1$.  \textbf{Converges absolutely}.

\subsection*{(c) \easy\ $\displaystyle\sum\frac{1}{n^{1+a}}$, $a>0$}

This is a $p$-series with $p=1+a$.  Converges $\Longleftrightarrow p>1\Longleftrightarrow a>0$. \textbf{Converges for all $a>0$}.

\subsection*{(d) \medium\ $\displaystyle\sum\frac{1}{n\sin(1/n)}$}

For small $u$, $\sin u\sim u$, so $n\sin(1/n)\to 1$ and $\dfrac{1}{n\sin(1/n)}\to 1\neq 0$.
By the \textbf{Necessary Condition}, the series \textbf{diverges}.

\subsection*{(e) \medium\ $\displaystyle\sum\frac{(1+1/n)^{n^3}}{n^3}$}

$(1+1/n)^n\to e$, so $(1+1/n)^{n^3}=\left[(1+1/n)^n\right]^{n^2}\sim e^{n^2}\to+\infty$.
Hence $a_n\to+\infty\neq 0$.  \textbf{Diverges}.

\subsection*{(f) \hard\ $\displaystyle\sum\frac{(2n)!}{4^n(n!)^2}$}

As in Exercise 2(i), by Stirling: $a_n\sim 1/\sqrt{\pi n}$.  \textbf{Limit Comparison} with $1/\sqrt{n}$ (which diverges, $p=1/2$) gives $L=1/\sqrt{\pi}\neq 0$, $\infty$.  \textbf{Diverges}.

\subsection*{(g) \hard\ $\displaystyle\sum\frac{1\cdot3\cdot5\cdots(2n-1)}{2\cdot4\cdot6\cdots(2n)}$}

\textbf{Ratio Test}: $\dfrac{a_{n+1}}{a_n}=\dfrac{2n+1}{2n+2}\to 1$.  Inconclusive.

By Wallis's product, $a_n=\dfrac{(2n-1)!!}{(2n)!!}\sim\dfrac{1}{\sqrt{\pi n}}$.  Since $\sum 1/\sqrt{\pi n}$ diverges, the series \textbf{diverges}.

\subsection*{(h) \hard\ $\displaystyle\sum\frac{(2n-1)!!}{(2n)!!}\cdot 3^{n/3}$}

Using $a_n\sim 1/\sqrt{\pi n}$ and multiplying by $3^{n/3}$, the terms grow like $3^{n/3}/\sqrt{\pi n}\to+\infty$.
\textbf{Diverges} by the Necessary Condition.

% -------------------------------------------------------
\section*{Exercise 4 — Theoretical Arguments}

\subsection*{(a) \medium\ $\sum a_n$ converges $\Rightarrow$ $\sum a_n/n^2$ converges absolutely}

Since $\sum a_n$ converges, $a_n\to 0$, so $|a_n|<1$ for all sufficiently large $n$.  Therefore for large $n$:
\[
  \left|\frac{a_n}{n^2}\right| = \frac{|a_n|}{n^2} \leq \frac{1}{n^2}.
\]
Since $\sum 1/n^2$ converges, by the \textbf{Comparison Test}, $\sum |a_n/n^2|$ converges.  $\square$

\subsection*{(b) \hard\ $\sum a_n$ converges $\Rightarrow$ $\sum \sin(a_n/n^3)$ converges absolutely}

Since $\sum a_n$ converges, $a_n\to 0$, so for large $n$, $|a_n/n^3|\leq |a_n|/n^3 \leq |a_n|$ which is eventually $< 1$, i.e.\ the argument of sine is small.  Using $|\sin u|\leq |u|$:
\[
  \left|\sin\!\left(\frac{a_n}{n^3}\right)\right| \leq \frac{|a_n|}{n^3} \leq \frac{|a_n|}{1} \cdot \frac{1}{n^3}.
\]
More carefully: for large $n$, $|a_n|\leq 1$, so $|a_n/n^3|\leq 1/n^3\leq 1$, giving $|\sin(a_n/n^3)|\leq |a_n|/n^3\leq |a_n|$.  Since $\sum|a_n|$ converges (actually we only know $\sum a_n$ converges), apply instead: $|\sin(a_n/n^3)|\leq 1/n^3$, and $\sum 1/n^3$ converges.  \textbf{Converges absolutely}. $\square$

\subsection*{(c) \medium\ Proof of Leibniz's Test}

Let $a_n\searrow 0$.  Define $S_N=\sum_{n=1}^N(-1)^{n+1}a_n$.

\textit{Even partial sums}: $S_{2N}=S_{2N-2}+(a_{2N-1}-a_{2N})\geq S_{2N-2}$ (non-decreasing), and $S_{2N}=a_1-(a_2-a_3)-\cdots-(a_{2N-2}-a_{2N-1})-a_{2N}\leq a_1$ (bounded above).

\textit{Odd partial sums}: $S_{2N+1}=S_{2N}+a_{2N+1}$ where $a_{2N+1}\to 0$, so $S_{2N+1}\to L$.

Both subsequences converge to the same limit $L$, so $S_N\to L$.

\textit{Error estimate}: $|S-S_N|\leq a_{N+1}$.  $\square$

\subsection*{(d) \medium\ Approximating $\sum\frac{(-1)^{n+1}}{n(n+1)^2}$ to within $10^{-5}$}

This is an alternating series with $a_n=1/(n(n+1)^2)$, decreasing to $0$.  By the \textbf{Leibniz error bound}, we need the first omitted term:
\[
  a_{N+1} = \frac{1}{(N+1)(N+2)^2} < 10^{-5}.
\]
Testing $N=46$: $a_{47}=\dfrac{1}{47\cdot 48^2}=\dfrac{1}{108288}\approx 9.2\times 10^{-6}<10^{-5}$.

Testing $N=45$: $a_{46}=\dfrac{1}{46\cdot47^2}=\dfrac{1}{101614}\approx 9.8\times 10^{-6}<10^{-5}$.

So summing the first $\mathbf{N=45}$ terms gives error $< 10^{-5}$.

\bigskip
\hrule
\bigskip
\textit{Key takeaways}: Alternating series are classified by checking absolute convergence first (Comparison/$p$-test), then conditional convergence via Leibniz.  The necessary condition eliminates many divergent cases quickly.  Stirling's approximation is essential for double-factorial and binomial-coefficient series.

\end{document}
